\section{Periodic-boundary ground space generated by a basis of normal tensors}
\label{sec:bnt_span_consequences}

Let $A_0,\ldots,A_{g-1}$ be irreducible canonical blocks and let every
weight $\mu_j$ be nonzero.  The preceding sections prove the simultaneous
word-span, one-step intersection, global boundary-closing, and periodic
chain-space statements at the source length of
\cite[Theorem~12]{PerezGarcia2007Matrix}.  They yield the exact kernel equality
under the source normalization
\begin{align}
    \sum_a A_j^a A_j^{a\,\dagger}&=\Id,
    &
    \sum_a A_j^{a\,\dagger}\Lambda_j A_j^a&=\Lambda_j,
    \qquad \Lambda_j>0.
    \notag
\end{align}
The first theorem below retains the earlier doubly-normalized specialization;
the second is the source-facing result.

\begin{theorem}[Doubly-normalized PGVWC07 source-length kernel equality]
    \label{thm:gs_eq_bnt_span}
    \lean{MPSTensor.ker_parentHamiltonian_toTensorFromBlocks_eq_bntMPSVectorSpan_of_global_cut_bnt_c1_pgvwc07}
    \leanok
    \uses{def:parent_hamiltonian, def:bnt_span, def:block_diagonal_tensor,
        def:irreducible_tensor, def:left_canonical_tensor,
        def:blocks_not_gauge_phase_equiv, def:l_blk_injective, def:mpv_overlap}
    % Scope restriction documented in
    % docs/paper-gaps/pgvwc07_ti_canonical_form_scope.tex.
    Let $d\ge1$ and $g\ge2$. For $0\le j<g$, let $D_j\ge1$ and
    let $A_j=(A_j^a)_{a=0}^{d-1}$ be an irreducible left-canonical block with
    normalized self-overlap. Assume that no two distinct blocks are equivalent
    up to gauge and phase. For every $0\le j<g$, let $\mu_j\ne0$. Suppose
    there is a common $L_0>0$ such that the words of length $L_0$ in each
    block $A_j$ span $\MN{D_j}$. In this doubly-normalized specialization,
    assume for every
    $0\le j<g$ that
    \begin{align}
        \sum_{a=0}^{d-1} A_j^a A_j^{a\,\dagger}&=\Id,
        \notag\\
        \sum_{a=0}^{d-1} A_j^{a\,\dagger}A_j^a&=\Id.
        \notag
    \end{align}
    The first identity is the normalization used in the source. In the general
    source canonical form, the second identity is replaced by a positive
    definite matrix $\Lambda_j$ satisfying
    \begin{align}
        \sum_{a=0}^{d-1} A_j^{a\,\dagger}\Lambda_j A_j^a&=\Lambda_j.
        \notag
    \end{align}
    Thus this theorem has the explicit specialization $\Lambda_j=\Id$ and is
    retained for compatibility; the unrestricted source result is stated below.
    If
    \begin{align}
        L&\ge 3(g-1)(L_0+1)+1,
        \notag\\
        N&\ge L+L_0,
        \notag
    \end{align}
    then the parent-Hamiltonian kernel is exactly the span of the BNT component
    vectors:
    \begin{align}
        \ker H_L^{(N)}\!\left(\bigoplus_j \mu_j A_j\right)
        &=
        \spn\{V^{(N)}(A_j):j=0,\ldots,g-1\}.
        \label{eq:parent_bnt_ground_space}
    \end{align}
\end{theorem}

\begin{proof}
    \leanok
    \uses{thm:gs_eq_bnt_span_source_normalization}
    % Scope restriction documented in
    % docs/paper-gaps/pgvwc07_ti_canonical_form_scope.tex.
    Apply the general-canonical-normalization theorem below with
    $\Lambda_j=\Id$.  Positive definiteness is immediate, and the assumed
    left-canonical identity is exactly the required dual fixed-point equation
    at $\Lambda_j=\Id$.
\end{proof}

\begin{theorem}[PGVWC07 source-length kernel equality at general canonical normalization]
    \label{thm:gs_eq_bnt_span_source_normalization}
    \lean{MPSTensor.ker_parentHamiltonian_toTensorFromBlocks_eq_bntMPSVectorSpan_of_global_cut_bnt_c1_pgvwc07_of_dualFixedPoint}
    \leanok
    \uses{def:parent_hamiltonian, def:bnt_span, def:block_diagonal_tensor,
        def:irreducible_tensor, def:blocks_not_gauge_phase_equiv,
        def:l_blk_injective}
    Retain the same nonzero dimensions and weights, irreducibility,
    gauge-phase inequivalence, common block-injectivity length,
    and source-length bounds, but do not assume left-canonicality. For every
    $0\le j<g$, assume there is a positive definite matrix $\Lambda_j$ such
    that
    \begin{align}
        \sum_{a=0}^{d-1} A_j^a A_j^{a\,\dagger}&=\Id,
        \notag\\
        \sum_{a=0}^{d-1} A_j^{a\,\dagger}\Lambda_j A_j^a&=\Lambda_j.
        \notag
    \end{align}
    Then
    \begin{align}
        \ker H_L^{(N)}\!\left(\bigoplus_j \mu_j A_j\right)
        &=
        \spn\{V^{(N)}(A_j):j=0,\ldots,g-1\}.
        \notag
    \end{align}
\end{theorem}

\begin{proof}
    \leanok
    \uses{cor:block_diagonal_kernel_global_cut,
        thm:chain_ground_space_eq_mpv_normal,
        lem:bnt_span_le_parent_hamiltonian_kernel}
    The source length bounds imply $L_0<L$, $L\le N$, and $L_0+1<N$.
    Theorem~\ref{thm:chain_ground_space_eq_mpv_normal} identifies each
    single-block periodic chain space with its matrix-product-vector line.
    Corollary~\ref{cor:block_diagonal_kernel_global_cut} therefore gives the
    containment of the parent-Hamiltonian kernel in the BNT span. Conversely,
    Lemma~\ref{lem:bnt_span_le_parent_hamiltonian_kernel} gives the reverse
    containment, and antisymmetry yields equality.
\end{proof}
